\chapter{Correcting What You Cannot Copy}\label{ch:qec}

\begin{etymology}
\textbf{Error} comes from the Latin \emph{errare}, ``to wander'', ``to stray''.
An error is a state that has wandered away from where it should be.

\textbf{Redundancy} comes from \emph{redundare}, ``to overflow'': \emph{re-},
``again'', and \emph{undare}, ``to surge like a wave''. More than is needed, kept
on purpose.
\end{etymology}

By the middle of 1994 quantum computing had an algorithm worth building a machine
for. Within weeks it also had an objection that seemed to make the machine
impossible. This chapter is about the objection and the answer, because the answer
is what turned the question ``can a quantum computer work?'' into the question this
book asks: ``how many qubits will it take?''.

\section{The objection}

\begin{history}{May 1994 \textperiodcentered{} Santa Fe Institute}
At a workshop at the Santa Fe Institute, a few weeks after Shor's result began to
circulate, Rolf Landauer of IBM argued that quantum computers could not correct
errors. Landauer was the physicist who had shown in 1961 that erasing information
costs energy, and he had spent decades studying what physics allows computers to
do. William Unruh published a similar warning about decoherence the next year.
\end{history}

Classical computers fight errors in three ways, and Shor himself later went through
them to show why the objection was so strong.

\begin{itemize}
\item \textbf{Checkpoints.} Save the state of the computation now and then; if an
error happens, go back to the last save. Saving the state is copying it.
\item \textbf{Error-correcting codes.} Add redundant parity bits that let you detect
and fix a flipped bit. The redundancy is, in effect, partial copies.
\item \textbf{Massive redundancy.} Run many copies, compare them, and take the
majority. If one of five copies fails, you rebuild the fifth from the other four.
That is copying too.
\end{itemize}

The no-cloning theorem (Chapter~\ref{ch:quantum}) forbids copying an unknown
quantum state. There is a second problem: classical errors are discrete (a bit
flips or it does not), while a qubit can drift by any small angle. And a third:
checking a qubit means measuring it, and measuring destroys the superposition.

Without error correction, every gate must be perfect. Factoring a number of
cryptographic size takes on the order of a billion steps, so each one would need an
error below one in a billion. Shor recalled that the best estimate at the time, by
Jeff Kimble, was that experimenters could do one or two quantum gates with about
80\% accuracy.

\section{Encoding without copying}

The simplest classical code repeats each bit three times: $0 \to 000$,
$1 \to 111$. If one bit flips, the majority still gives the right answer. Asher
Peres had already noted in 1985 that the same thing can be written for qubits:
\[
  \alpha\ket0 + \beta\ket1 \;\longrightarrow\; \alpha\ket{000} + \beta\ket{111}.
\]
This is not cloning. Cloning would produce
$(\alpha\ket0 + \beta\ket1)^{\otimes3}$, three independent copies of the
superposition. The encoded state is one entangled superposition spread over three
qubits.

The trick that gets around measurement is to ask only about parities. ``Are qubits 1
and 2 equal? Are qubits 2 and 3 equal?'' Those questions can be answered with extra
qubits without learning anything about $\alpha$ or $\beta$. The answers, called the
\emph{syndrome}, say which qubit flipped, if any, and a gate flips it back. The
superposition is never measured.

And the problem of continuous errors dissolves on its own. A small drift is a
superposition of ``no error'' and ``a flip''. Measuring the syndrome forces it to
choose: either nothing happened, or a full flip happened, which can then be fixed.
Measurement, the enemy of superposition, turns analog errors into digital ones.

\section{Shor's nine qubits}

The three-qubit code fixes bit flips, but it does nothing for phase errors, the
error that turns $\ket0 + \ket1$ into $\ket0 - \ket1$. Worse, a phase error on any of
the three qubits becomes a phase error on the encoded qubit, so the code makes phase
errors three times more likely.

\begin{history}{1995 \textperiodcentered{} Bell Labs}
Shor noticed that a phase error is a bit flip seen in another basis: apply a
Hadamard and $Z$ becomes $X$ (Chapter~\ref{ch:state}). So he protected against phase
errors with a three-qubit code in the Hadamard basis, and protected each of those
three qubits against bit flips with another three-qubit code. Three times three: one
logical qubit in nine physical ones, protected against any error on a single qubit.
The paper appeared in 1995. In 1996, Robert Calderbank and Shor, and independently
Andrew Steane at Oxford, showed how to turn pairs of classical codes into quantum
ones. Steane's code uses seven qubits. The family is called CSS, after the three of
them.
\end{history}

\begin{concept}{Logical qubit}
A qubit stored in a collective state of many physical qubits, so that errors on a
few of them can be detected and corrected without measuring the information itself.

\textbf{What it fixes:} the computation no longer needs perfect hardware, only
hardware that is good enough.

\textbf{What it costs:} many physical qubits per logical qubit, constant syndrome
measurements, and fast classical processing to decode them, for the whole duration
of the computation.
\end{concept}

\section{From possible to engineering}

A code that fixes one error is not enough. The extra qubits and gates of the code
can fail too, and if the code introduces more errors than it removes, it makes
things worse.

\begin{history}{1996--1997 \textperiodcentered{} the threshold theorem}
Several groups proved, independently, the result that settled the question: Dorit
Aharonov and Michael Ben-Or; Emanuel Knill, Raymond Laflamme and Wojciech Zurek;
Alexei Kitaev; and building on Shor's own 1996 work on fault-tolerant gates. If the
error rate of every physical operation is below a certain \emph{threshold}, codes
can be nested or enlarged so that the logical error rate falls as low as you want,
with an overhead that grows only polylogarithmically with the length of the
computation.
\end{history}

The threshold theorem is the reason quantum computing is an engineering problem and
not a dead end. It does not say the engineering is easy. It says how to measure the
distance: get the physical error rate below the threshold, then pay in qubits for
every factor of ten in reliability.

In 1997 Kitaev proposed codes laid out on a lattice, first on a torus and then, with
Sergey Bravyi in 1998, on a flat surface. Every check involves only neighboring
qubits, which suits chips where qubits can only talk to their neighbors.
These \emph{surface codes} have a threshold of about one percent, high by the
standards of the field, and became the default plan for superconducting hardware.
Their price is overhead: hundreds to thousands of physical qubits per logical qubit,
depending on how good the physical qubits are. That overhead is why the estimate for
breaking RSA-2048 in Chapter~\ref{ch:factoring} is counted in millions of physical qubits rather than
the few thousand logical qubits the algorithm needs.
